%------------------------------------------------------------------------------%
\section{Teorema Fundamental do Cálculo}
\label{sec:teorema_fundamental_calculo}

Nessa seção veremos um resultado importante e que facilita o cálculo das integrais
definidas: O Teorema Fundamental do Cálculo. Esse Teorema será dividido as duas
partes e estabelece uma conexão entre dois tópicos importantes do cálculo:
integração e derivação.

Antes de falarmos desse importante Teorema precisamos de um resultado auxiliar
chamado: Teorema do valor médio para integrais.

No início do capítulo anterior, definimos como encontrar a área sob uma curva
$y=f(x)$ contínua e não-negativa no intervalo $[a, b]$, conforme mostra a
Figura~\ref{fig:areasobumacurva2}.

%----------------------------------------------------------------------------------
\begin{figure}[h!]
   \centering
   \input{fig_integral_definida}
   \caption{Área sob uma curva.}
   \label{fig:areasobumacurva2}
\end{figure}
%----------------------------------------------------------------------------------

A pergunta que surge agora é: Será que existe $c \in [a, b]$, de modo que a área
sob a curva $y=f(x)$ seja igual área do retângulo de altura $f(c)$ e base $b-a$?
O Teorema do Valor médio para integrais responde esse questionamento.

%----------------------------------------------------------------------------------
%\begin{figure}
%   \centering
%   \input{fig_valor_medio_minimo}
%   \caption{Teorema valor médio -- mínimo.}
%   \label{fig:retangulos}
%\end{figure}
%----------------------------------------------------------------------------------

%----------------------------------------------------------------------------------
%\begin{figure}
%   \centering
%   % Title: Figure 1
% Creator: GL2PS 1.4.2, (C) 1999-2020 C. Geuzaine
% For: Octave
% CreationDate: Fri Aug  4 13:35:06 2023
\setlength{\unitlength}{1pt}
\begin{picture}(0,0)
\includegraphics[width=14cm,height=5cm]{fig_valor_medio_maximo}
\end{picture}%
\begin{picture}(396,141)(0,0)
\fontsize{10}{0}\selectfont\put(92.9927,20.5695){\makebox(0,0)[t]{\textcolor[rgb]{0.15,0.15,0.15}{{$a$}}}}
\fontsize{10}{0}\selectfont\put(300.004,20.5695){\makebox(0,0)[t]{\textcolor[rgb]{0.15,0.15,0.15}{{$\;b$}}}}
\fontsize{10}{0}\selectfont\put(113.798,20.5695){\makebox(0,0)[t]{\textcolor[rgb]{0.15,0.15,0.15}{{$u$}}}}
\fontsize{10}{0}\selectfont\put(241.749,20.5695){\makebox(0,0)[t]{\textcolor[rgb]{0.15,0.15,0.15}{{$v$}}}}
\fontsize{10}{0}\selectfont\put(58.4269,121.068){\makebox(0,0)[r]{\textcolor[rgb]{0.15,0.15,0.15}{{$f(u)$}}}}
\fontsize{10}{0}\selectfont\put(58.4269,51.8252){\makebox(0,0)[r]{\textcolor[rgb]{0.15,0.15,0.15}{{$f(v)$}}}}
\fontsize{10}{0}\selectfont\put(311.833,110.104){\makebox(0,0)[l]{\textcolor[rgb]{0,0,0}{{$y=f(x)$}}}}
\end{picture}

%   \caption{Teorema valor médio -- máximo.}
 %  \label{fig:retangulos}
%\end{figure}
%----------------------------------------------------------------------------------

\begin{teorema}{Teorema do Valor Médio para Integrais}{teo:tvmi}
  Seja $f$ uma função contínua no intervalo $[a, b]$, então existe $c \in (a, b)$
  tal que
  \[
      f(c)(b-a)= \int_a^b f(x)dx
  \]
  ou equivalentemente,
  \[
      f(c)= \frac{1}{b-a}\int_a^b f(x)dx.
  \]
  O valor $f(c)$ é chamado \textit{Valor médio} de $f$ no intervalo $[a, b]$
  e denotamos por $f_{\text{médio}}$.
\end{teorema}

Geometricamente, o Teorema \ref{teo:tvmi} afirma que existe um $c \in (a, b)$
tal que
\[
  \int_a^bf(x) dx
\]
(a área sob o gráfico de f) é igual a
$f(c) (b-a)$ (a área do retângulo de base $[a, b]$ e altura $f(c)$),
como ilustra a Figura~\ref{fig:valormedio_media}.

\begin{figure}
   \centering
   \input{fig_valor_medio_media}
   \caption{Teorema valor médio -- média.}
   \label{fig:valormedio_media}
\end{figure}

\begin{proof}
  Para iniciarmos a demonstração, devemos lembrar que como $f$ é contínua em
  $[a, b]$, então $f$ assume um
  valor mínimo $f(u)$ e um
  valor máximo $f(v)$ com $u, v \in [a, b]$
  (Teorema de Weierstrass~\ref{teo:weistress}).
  Assim,
  \[
    f(u) \leq f(x) \leq f(v)
  \]
  Como já vimos que a integral preserva desigualdade, temos que
  \[
    \int_a^b f(u) dx \leq \int_a^b f(x) dx\leq \int_a^b f(v) dx
  \]
  Portanto,
  \[
    f(u)(b-a) \leq \int_a^b f(x) dx \leq f(v)(b-a)
  \]
  Dividindo essa última desigualdade por $b-a$, chegamos que
  \[
    f(u) \leq \frac{1}{b-a}\int_a^b f(x)dx \leq f(v)
  \]
  Como $f$ é contínua em $[a, b]$, pelo Teorema do Valor
  Intermediário~\ref{teo:valor-intermediario}, $f$ deve assumir
  todos os valores entre o mínimo, $f(u)$, e o máximo, $f(v)$,
  de $f$. Portanto, deve assumir o valor
  \[
    \frac{1}{b-a}\int_a^b f(x)dx
  \]
  para algum valor $c \in (a, b)$. Em outras palavras, existe
  $c \in (a, b)$ tal que
  \[
    f(c)= \frac{1}{b-a}\int_a^b f(x)dx \,.
  \]
\end{proof}

Observe que a continuidade de $f$ em $[a, b]$ foi uma hipótese crucial na
demonstração do Teorema~\ref{teo:tvmi}. Sem essa hipótese, é possível que
uma função descontínua não assuma seu valor médio.

%------------------------------------------------------------------------------%
\begin{teorema}{Teorema Fundamental do Cálculo -- Parte 1}{teo:fundamental-1}
  Se $f$ for contínua em $[a, b]$, então a função $F$ definida por
  \[
    F(x)= \int_a^x f(t)dt, \hspace{1cm} a \leq x \leq b
  \]
  é contínua em $[a, b]$ e derivável em $(a, b)$ e $F'(x)=f(x)$
  note que em outras palavras temos que
  \[
    \frac{d}{dx} \left( \int_a^x f(t)dt \right)=f(x)
  \]
\end{teorema}

%------------------------------------------------------------------------------%
\begin{example}
  Calcule a derivada da função \(g(x)=\int_0^x \sqrt{1+t^2} dt\)
  \solution
  Sendo $f(t)= \sqrt{1+t^2}$ contínua em todo o conjunto dos números reais,
  pela parte 1 do Teorema Fundamental do Cálculo
  \[
    g'(x)=\sqrt{1+x^2}
  \]
\end{example}
%------------------------------------------------------------------------------%

%------------------------------------------------------------------------------%
\begin{example}
  Calcule a derivada da função \(h(x)=\int_1^{x^4} \sec(t)dt\)
  \solution
  Nesse item devemos ter cuidado ao aplicarmos o Teorema Fundamental do
  Cálculo uma vez que $h(x)=p(q(x))$, onde
  \[
    p(x)=\int_0^x \sec(t)dt
    \qquad\text{e}\qquad
    q(x)= x^4
  \]
  Portanto, para encontrar
  \[
    \frac{d}{dx}\big(h(x)\big)
  \]
  temos que usar a regra da
  cadeia. Ora, o Teorema Fundamental do Cálculo Parte 1 nos diz que
  \[
    \frac{d}{dx}\big(p(x)\big)= \frac{d}{dx} \int_0^x \sec(t)dt= \sec(x),
  \]
  além disso,
  \[
  \frac{d}{dx}\left(x^4\right) = 4x^3
  \]
  Assim, pela regra da cadeia
  \[
    \frac{d}{dx}\int_1^{x^4} \sec(t)dt= 4x^3\sec\left(x^4\right)
  \]
\end{example}
%------------------------------------------------------------------------------%

%------------------------------------------------------------------------------%
\begin{example}
  Calcule a derivada da função \(p(x)= \int_{1-2x}^{1+2x} t\sin(t)dt\)
  \solution
  Inicialmente, vamos usar as propriedades de integral para manipular a
  integral. Pela propriedade 7. de integral definida, podemos escrever
  \[
      \int_{1-2x}^{1+2x} t\sin(t)dt
      = \int_{1-2x}^{0} t\sin(t)dt
      = \int_{0}^{1+2x} t\sin(t)dt
  \]
  Pela propriedade 1. de Integral definida
  \[
      \int_{1-2x}^{1+2x} t\sin(t)dt
      = - \int_{0}^{1-2x} t\sin(t)dt
      + \int_{0}^{1+2x} t\sin(t)dt
  \]
  Usando a parte 1 do Teorema Fundamental do Cálculo e a Regra da cadeia,
  temos
  \begin{align*}
    p'(x)
    & = \frac{d}{dx} \int_{1-2x}^{1+2x} t\sin(t)dt \\
    & = \frac{d}{dx} \left(
    - \int_{0}^{1-2x} t\sin(t)dt
    + \int_{0}^{1+2x} t\sin(t)dt
    \right) \\
    & = 2x \sin(x)+2x \sin(x) \\
    & = 4x\sin(x).
  \end{align*}

\end{example}
%------------------------------------------------------------------------------%

%------------------------------------------------------------------------------%
\begin{example}
  Determine
  \(f'\left(\frac{\pi}{4}\right)\), onde
  \(f(x)= e^{h(x)}\) e
  \(h(x)= \int^{\tan(x)}_1 e^t dt\)
  \solution
  Note que $f$ é uma função composta, para derivar usamos a regra da cadeia.
  Assim,
  \[
    f'(x)= e^{h(x)}h'(x)
  \]
  Para encontrar $h'(x)$ usamos o Teorema Fundamental do Cálculo parte 1,
  combinado com a regra da cadeia. Portanto,
  \[
    h'(x)=e^{\tan(x)}\sec^2(x)
  \]
  Logo,
  \[
    f'(x)
    = e^{h(x)}h'(x)
    =e^{h(x)}e^{\tan(x)}\sec^2(x)
  \]
  e
  \[
    f'\left(\frac{\pi}{4}\right)
    = e^{h(\sfrac{\pi}{4})}
      e^{\tan(\sfrac{\pi}{4})}
      \sec^2\left(\frac{\pi}{4}\right)
    = 0 \cdot e^1 \cdot \left(\frac{2}{\sqrt{2}}\right)^2=0
  \]
  onde,
  \[
    h\left(\frac{\pi}{4} \right)= \int_1^1 e^t dt=0.
  \]
\end{example}
%------------------------------------------------------------------------------%

%------------------------------------------------------------------------------%
\begin{teorema}{O Teorema Fundamental do Cálculo -- Parte 2}{teo:fundamental-2}
  Se $f$ for contínua em $[a, b]$ então
  \[
    \int_a^{b}f(x)dx= F(b)-F(a)
  \]
  onde $F$ é qualquer primitiva de $f$, isto é $F'=f$
\end{teorema}
%------------------------------------------------------------------------------%

%------------------------------------------------------------------------------%
\begin{proof}
  Seja
  \[
    G(x)= \int_a^x f(x)dx \,.
  \]
  Pela parte 1 do Teorema Fundamental do Cálculo,
  $g$ é uma primitiva de $f$. Se $F$ for qualquer outra primitiva de $f$,
  então sabemos que $F$ e $G$ se diferem por uma constante, ou seja:
  \[
      F(x)= G(x)+K
  \]
  Assim,
  \begin{align*}
    F(b)-F(a)
      & = \big[ G(b) + K \big]
        - \big[ G(a) + K \big] \\
      & = \int_a^b f(x)dx -  \int_a^a f(x)dx \\
      & = \int_a^b f(x)dx \,.
  \end{align*}
\end{proof}
%------------------------------------------------------------------------------%

O Teorema \ref{teo:fundamental-2} nos diz que para calcular a integral
definida
\[\int_a^b f(x)dx,\]
precisamos fazer duas coisas:
\begin{enumerate}
    \item Determinar uma primitiva $F$ qualquer para $f$.
    \item Calcular o número $F(b)-F(a)$.
\end{enumerate}
Além disso, vale ressaltar que enquanto a integral indefinida
\[\int f(x)dx\]
é uma função (ou uma família de funções),
a integral definida
\[\int_a^b f(x)dx\]
é um número.

%------------------------------------------------------------------------------%
Agora temos uma ferramenta importante para calcular integrais definidas.
Vamos aos exemplos.
\begin{example}
  Calcule \(\int_1^3 e^x dx\)
  \solution
  Inicialmente, note que $f(x)=e^x$ é contínua no intervalo $[1, 3]$. Além
  disso, a primitiva mais geral para $f$ é $G(x)=e^x+ C$, como no Teorema
  Fundamental do cálculo precisamos de qualquer primitiva, então pegamos a
  primitiva mais simples $F(x)=e^x$. Logo, usando o Teorema Fundamental do
  Cálculo 2 temos
  \[
    \int_1^3 e^x dx
    = e^x \evalat{1}{3}
    = e^3-e^1
    = e^3 - e.
  \]
\end{example}
%------------------------------------------------------------------------------%

%------------------------------------------------------------------------------%
\begin{example}
  Calcule \(\int_3^6 \frac{dx}{x}\)
  \solution
  Note que $f(x)=\sfrac{1}{x}$ é contínua no intervalo $[3, 6]$. Além disso, a
  primitiva mais geral para $f$ é $G(x)=\ln\abs{x}+ C$, como no Teorema Fundamental
  do cálculo precisamos de qualquer primitiva, então pegamos a primitiva mais
  simples $F(x)=\ln\abs{x}$. Logo, usando o Teorema Fundamental do Cálculo 2 temos
  \[
    \int_3^6 \frac{dx}{x}
    = \ln\abs{x} \evalat{1}{3}
    = \ln(6) - \ln(3)
    = \ln\left(\frac{6}{3}\right)
    = \ln(2).
  \]
\end{example}
%------------------------------------------------------------------------------%

%------------------------------------------------------------------------------%
\begin{example}
  Calcule
  \(
    \int_0^1 \left[ 20x^5-e^x+ 3x^6- \cos x+ \frac{3}{1+x^2}-3 \right] dx
  \)
  \solution
  Note que
  \[
    f(x)=20x^5-e^x+ 3x^6- \cos x+ \frac{3}{1+x^2}-3
  \]
  é contínua no
  intervalo $[0, 1]$, uma vez que é soma de funções contínuas.

  Para calcularmos a integral definida pedida, iniciaremos primeiro calculando
  a integral indefinida, usando a linearidade da integral.
  \begin{align*}
    F(x)
    & = \int \left[ 20x^5-e^x+ 3x^6 - \cos x+  \frac{3}{1+x^2}-3 \right] dx \\
    & = \int \left[ 20x^5-e^x+ 3x^6 - \cos x+ 3\frac{1}{1+x^2}-3 \right] dx \\
    & = 20 \int x^5 dx
      -    \int e^x dx
      + 3  \int x^6 dx
      -    \int \cos(x) dx \\
    & \hspace{10mm}
      + 3  \int \frac{1}{1+x^2} dx
      - 3  \int dx \\
    & = \frac{20x^6}{6}
      - e^x
      + \frac{3x^7}{7}
      - \sin(x)
      + 3\arctan(x)
      - 3x + C
  \end{align*}
  Para encontrarmos a integral definida, usamos o Teorema Fundamental do
  Cálculo parte 2 temos que
  \begin{align*}
    I
    & = \int_0^1 \left[ 20x^5-e^x+ 3x^6- \cos x+ \frac{3}{1+x^2}-3 \right] dx \\
    & = \left[ \frac{20x^6}{6}
      - e^x
      + \frac{3x^7}{7}
      - \sin(x)
      + 3\arctan(x)
      - 3x \right] \evalat{0}{1} \\
    & = \left( \frac{20}{6}
      - e
      + \frac{3}{7}
      - \sin(1)
      + 3 \arctan(1)
      - 3 \right) \\
    & \hspace{10mm}
      - \left( 0 - e^0 + 0 - \sin(0) + 3\arctan(0) - 0 \right) \\
    & = \left(\frac{10}{3}-e+\frac{3}{7}-\sin(1) + \frac{3\pi}{4} - 3 \right) - (-1) \\
    & = \frac{10}{3} - e + \frac{3}{7} - \sin(1) + \frac{3\pi}{4} - 3 + 1 \\
    & = \frac{148 - 84e - 84 \sin(1) + 63 \pi}{84}
  \end{align*}
\end{example}
%------------------------------------------------------------------------------%


%------------------------------------------------------------------------------%
\begin{example}
  Encontre a área sob a parábola $y=x^2$ de 0 até 1.
  \solution
  A figura a seguir mostra o gráfico da função $y=x^2$.

  \IncludeGraphicCentred{exemplo_area_sob_parabola}[0.5][-1][-1.5]

  Note que $y=x^2$ é uma função positiva no intervalo $[0, 1]$
  então
  \[
    A
    = \int^{0}_{1} x^2 dx
    = \frac{x^3}{3} \evalat{0}{1}
    = \frac{1}{3}-0
    = \frac{1}{3}
  \]
  onde, para calcularmos a integral definida usamos o
  Teorema Fundamental do Cálculo parte 2.
\end{example}
%------------------------------------------------------------------------------%

%------------------------------------------------------------------------------%
\begin{example}
  Encontre a área sob a curva $y=\cos(x) $ de $0$ até $b$, onde
  \(0 \leq b \leq \sfrac{\pi}{2}\)
  \solution
  Note que a função $y=\cos(x)$ é positiva e contínua no intervalo $[0, b]$,
  uma vez que \(b \in \left[0, \sfrac{\pi}{2}\right]\), conforme a figura .

  \IncludeGraphicCentred{exemplo_area_sob_cosseno}[0.5][-1.5][-1.5]

  Portanto, usando o Teorema Fundamental do Cálculo parte 2
  \[
    A
    = \int_0^b \cos(x) dx
    = \sin(x) \evalat{0}{b}
    = \sin(b) - \sin(0)
    = \sin(b) \,.
  \]
\end{example}
%------------------------------------------------------------------------------%

%------------------------------------------------------------------------------%
\begin{example}
  O que está errado no seguinte cálculo
  \[
    \int_{-1}^3 \frac{1}{x^2} dx
    = \frac{x^{-1}}{-1} \evalat{-1}{3}
    = - \frac{1}{3} - 1
    = - \frac{4}{3}
  \]
  \solution
  Para começarmos podemos observar, pela propriedade 8 de integrais, que esse
  resultado está errado, uma vez que
  \[
    f(x) = \frac{1}{x^2} \geq 0
  \]
  no intervalo
  $[-1, 3]$ teríamos que ter que
  \[
    \int_{-1}^3 \frac{1}{x^2} dx \geq 0.
  \]
  Portanto, devemos estar atentos as hipóteses quando quisermos usar o Teorema
  fundamental do cálculo. Note que
  \[
    f(x)=\frac{1}{x^2}
  \]
  não é contínua no intervalo $[-1, 3]$, uma vez que existe uma descontinuidade
  em $x=0$ ($f$ não está definida em $x=0$).
  Nesse caso, não podemos usar cegamente o Teorema
  Fundamental do Cálculo. Veremos como resolver esse tipo de problema quando
  estudarmos \emph{Integrais impróprias}.
\end{example}
%------------------------------------------------------------------------------%

%------------------------------------------------------------------------------%
\begin{example}
Seja $f(x)= x^2+2$
\begin{enumerate}[label=\alph*)]
  \item Encontre o valor médio (média) de $f(x)$ no intervalo $[-1, 3]$.
  \item Encontre o(s) valor(es) de $c$ que satisfaz(em)
        o Teorema do Valor médio para integrais.
\end{enumerate}
\solution
\begin{enumerate}[label=\alph*)]
  \item Observando que $f(x)$ é contínua $[-1, 3]$, uma vez que é um polinômio.
    Então o valor médio de $f$ no intervalo $[-1, 3]$ é:
    \begin{align*}
      f_{\textrm{médio}}
      & = \frac{1}{b-a}\int_{-1}^{3} f(x)dx \\
      & = \frac{1}{3-(-1)}\int_{-1}^{3} (x^2+2)dx
    \end{align*}
    Usamos agora o Teorema Fundamental do Cálculo
    \begin{align*}
      f_{\textrm{médio}}
      & = \frac{1}{4}\left[ \frac{x^3}{3}+2x\right] \evalat{-1}{3} \\
      & = \frac{1}{4}\left( \frac{27}{3}+6\right)
        - \frac{1}{4}\left( \frac{-1}{3}-2\right) \\
      & = \frac{13}{3}
    \end{align*}

  \item O(s) valor(es) de $c$ que atende(m) o Teorema do valor médio para
    integrais \ref{teo:tvmi} satisfaz(em)
    \[
      f(c) = f_{\textrm{médio}}= \frac{13}{3}.
    \]
    Uma vez que $f(c)= c^2+2$, temos
    \begin{align*}
      c^2+2 & = \frac{13}{3}                 \\
      c^2   & = \frac{13}{3}-2 = \frac{7}{3} \\
      c     & = \pm \sqrt{\frac{7}{3}}
    \end{align*}
    Assim existem dois valores para $c$
    \[
      c= \pm \sqrt{\frac{7}{3}} \approx \pm 1,5275
    \]
    no intervalo $[-1, 3]$, que satisfazem o Teorema do Valor Médio para integrais.
\end{enumerate}
\end{example}
%----------------------------------------------------------------------------------

Podemos justapor as duas partes do Teorema Fundamental do cálculo.

\begin{teorema}{Teorema Fundamental do Cálculo}{teo:fundamental}
  Se $f$ for contínua em $[a, b]$ valem as afirmações:
  \begin{enumerate}
    \item Se \(a \leq x \leq b\)
      \[
        F(x)  = \int_a^x f(t)dt
        \qquad\Rightarrow\qquad
        F'(x) = f(x)
      \]
    \item Se $F$ for qualquer primitiva de $f$, isto é $F'=f$, então
      \[
        \int_a^{b}f(x)dx= F(b)-F(a)
      \]
  \end{enumerate}
\end{teorema}

%------------------------------------------------------------------------------%
